\chapter{工程模板与评审清单}
\label{chap:engineering-templates}

工程文档的价值不在篇幅，而在于能否减少歧义、驱动验证并支持变更。模板应按项目风险裁剪：小型实验可以使用一页 Markdown，大型产品可以把同样字段分布到需求、接口、测试和发布系统中。字段可以简化，但关键决策不能只存在于口头沟通中。

\section{需求条目模板}

一条可执行需求至少包含：
\begin{description}
  \item[ID] 稳定且唯一，例如 \texttt{SYS-BOOT-017}；
  \item[陈述] 在明确条件下，单一主体必须满足的可观察行为；
  \item[理由] 用户目标、法规、安全目标或上游约束；
  \item[适用条件] 硬件版本、工作模式、环境和输入范围；
  \item[验收标准] 数值范围、统计口径、允许误差和失败判据；
  \item[验证方法] 分析、检查、测试或演示，以及对应测试 ID；
  \item[追踪关系] 上游需求、架构元素、接口、实现和发布证据；
  \item[状态] 草拟、批准、实现、验证或废弃，以及责任人。
\end{description}

示例：
\begin{quote}
在输入电压 12\,V、环境温度 $-20\sim 60\,^{\circ}\mathrm{C}$、已安装有效生产镜像且文件系统完成正常关闭的条件下，设备从上电到核心采集服务报告 ready 的第 99.9 百分位时间不得超过 8\,s。测试应覆盖不少于 1000 次冷启动，并保存每次启动阶段时间戳。
\end{quote}

该写法仍需由项目定义电压容差、温度驻留、ready 的语义和异常样本处理方式，但已经能够驱动启动埋点、测试工装和统计报告。

\section{接口契约模板}

接口可以是函数、消息、共享内存、寄存器、连接器或网络协议。建议按以下结构记录：

\begin{table}[htbp]
  \centering
  \caption{接口契约字段}
  \begin{tabular}{p{30mm}p{110mm}}
    \hline
    字段 & 内容 \\
    \hline
    标识与版本 & 稳定 ID、提供方、使用方、版本协商和弃用策略 \\
    数据语义 & 类型、单位、范围、分辨率、字节序、对齐、保留位和无效值 \\
    时间语义 & 周期、截止期、超时起点、时间戳时钟域、抖动和新鲜度 \\
    所有权 & 缓冲区、内存、电源、时钟、总线和设备状态由谁管理 \\
    并发语义 & 可重入性、调用上下文、锁、阻塞、队列深度和背压 \\
    故障语义 & 错误码、短传输、重复、乱序、丢失、重试和部分成功 \\
    生命周期 & 上电默认值、初始化、暂停、复位、升级、热插拔和关闭 \\
    安全属性 & 身份、权限、完整性、机密性、防重放和审计要求 \\
    验证资产 & schema、静态断言、兼容测试、模糊测试和抓包/trace 方法 \\
    \hline
  \end{tabular}
\end{table}

对于跨核或持久化接口，还应明确 magic、长度、schema 版本、generation、校验和及提交标志；对于电气接口，还应补充电平、驱动能力、上拉、保护、连接顺序和故障电流。

\section{资源预算表}

预算必须同时记录目标、当前测量、最坏条件、裕量和责任人。只写“优化 CPU”无法判断是否完成。

\begin{table}[htbp]
  \centering
  \caption{资源预算示例}
  \begin{tabular}{p{24mm}p{25mm}p{40mm}p{38mm}}
    \hline
    资源 & 系统上限 & 模块分配与当前证据 & 裕量/动作 \\
    \hline
    启动时间 & 8\,s P99.9 & 采集服务 1.2\,s，冷启动 trace & 保留 20\% \\
    常驻内存 & 256\,MiB & 服务 36\,MiB，峰值 RSS & 修复增长趋势 \\
    DDR 带宽 & 4.0\,GB/s & 视频 2.2\,GB/s，总线计数器 & 压测并发显示 \\
    功耗 & 8\,W & SoC 3.5\,W，高温稳态 & 校核散热裕量 \\
    写入寿命 & 5 年 & 日志 40\,MiB/日，写放大实测 & 限流与轮转 \\
    \hline
  \end{tabular}
\end{table}

测量必须说明工具、采样窗口和负载。平均值不能替代峰值，实验室典型温度也不能替代产品规定的最坏环境。

\section{架构决策记录}

ADR 适合记录难以从最终代码反推出理由的决定，推荐字段为：
\begin{enumerate}
  \item 标题、日期、状态和决策人；
  \item 背景、必须满足的约束和不在范围内的问题；
  \item 至少两个可行选项及其证据；
  \item 决定和直接理由；
  \item 正面、负面和操作层面的后果；
  \item 重新评估的触发条件；
  \item 关联需求、原型、测量和后续任务。
\end{enumerate}

ADR 不应被改写成“当时就应该知道”的完美历史。若条件变化，应新增一条替代决策并链接旧记录，从而保留推理链。

\section{风险与故障假设}

风险条目应区分事实、假设和未知项。最小字段包括触发条件、概率、影响、可检测性、缓解、应急措施、责任人和关闭证据。

对安全、可靠性或实时性关键功能，还应维护故障假设：单点故障还是多点故障、持续还是瞬态、是否允许共同原因、检测时间和进入安全状态的最大时间。测试若没有覆盖设计所依赖的故障假设，不能证明恢复机制有效。

\section{测试用例模板}

测试用例建议包含：
\begin{itemize}
  \item 测试 ID、关联需求和目标风险；
  \item DUT 硬件版本、软件 manifest、配置和校准状态；
  \item 仪器、拓扑、环境和前置状态；
  \item 输入序列、时序和随机种子；
  \item 可机器判断的通过/失败条件；
  \item 原始证据位置和摘要生成方法；
  \item 清理、恢复和重复执行要求。
\end{itemize}

一组完整测试不只覆盖正常值，还应覆盖边界值、无效输入、资源耗尽、超时、重复、乱序、并发、复位、掉电和版本不兼容。涉及概率性指标时，应预先定义样本量、置信度和异常值处理规则，避免看到结果后改变统计口径。

\section{缺陷报告与根因分析}

高质量缺陷报告应让另一位工程师能够复现并判断影响范围：
\begin{itemize}
  \item 首次出现版本、最近已知正常版本和复现概率；
  \item 最小复现步骤、实际结果、期望结果和原始日志；
  \item 电源、温度、网络、外设和操作历史；
  \item 故障时各处理器、进程、线程和设备状态；
  \item 临时绕过方法及其副作用；
  \item 严重度、用户影响和是否可能造成数据或安全损失。
\end{itemize}

根因分析应区分故障触发、缺陷根因、逃逸原因和放大因素。修复不仅要改变触发位置的代码，还要增加能够防止同类问题再次逃逸的测试、检查或设计约束。

\section{发布清单}

发布评审至少确认：
\begin{itemize}
  \item 需求覆盖、已知缺陷和残余风险已经批准；
  \item 源码、配置、工具链、依赖和构建环境可以复现；
  \item 镜像、符号、SBOM、许可证清单和摘要已归档；
  \item 签名、版本、防回滚、兼容矩阵和密钥角色正确；
  \item 安装、升级、降级限制、断电恢复和出厂恢复已经验证；
  \item 生产工装、校准、设备身份和追溯数据版本匹配；
  \item 日志、指标、诊断包和现场支持手册能够识别该版本；
  \item 分批发布、健康门槛、停止条件和回滚责任人明确。
\end{itemize}

发布完成后，应从独立存档重建或安装一次，而不是只验证开发机工作区中的产物。

\section{变更影响检查}

评审任何跨层变更时，依次检查：
\begin{enumerate}
  \item 需求和外部可见行为是否变化；
  \item 引脚、寄存器、Device Tree、消息、文件格式或网络 ABI 是否变化；
  \item 启动、实时、内存、带宽、功耗、热和寿命预算是否变化；
  \item 安全边界、密钥、权限、功能安全假设是否变化；
  \item 制造、校准、维修、升级、回滚和数据迁移是否变化；
  \item 哪些测试证据失效，必须在哪些硬件版本上重新执行。
\end{enumerate}

如果评审结论是“无影响”，仍应记录分析过的边界。明确的无影响证据比没有提及更容易在后续审计和故障调查中复用。

\section{模板使用原则}

模板是思考提示，不是合规装饰。项目早期允许字段未知，但必须标出责任人和关闭时点；项目后期允许保留风险，但必须说明接受者、影响和监测措施。任何模板一旦长期出现无人阅读、无法验证或与实现重复维护的字段，都应合并、自动生成或删除。
